% =========================================================
% Subsequences & Bolzano–Weierstrass — Curated Exercise Set
% Sources: Abbott, Bartle & Sherbert, Bruckner, Tao, Lebl
% =========================================================

\section*{Subsequences and Bolzano–Weierstrass — Exercise Review Set}

These exercises develop the central extraction techniques of real analysis.
Subsequences are the bridge between boundedness and compactness, and they
lead to the Bolzano–Weierstrass theorem.

\bigskip

% ---------------------------------------------------------
% Stub Macro
% ---------------------------------------------------------

\newcommand{\exstub}[4]{
\subsection*{Exercise: #1}

\textbf{Source:} #2

\textbf{Technique:} #3

\textbf{Task.} #4

\begin{remark*}[Work Area]
Write your proof or computation here.
\end{remark*}

\bigskip
}

% =========================================================
% ABBOTT
% =========================================================

\section*{Exercises from Abbott}

\exstub
{Abbott 2.4.1}
{Abbott, \textit{Understanding Analysis}}
{unpack-definition}
{Define a subsequence of a sequence and give two examples.}

\exstub
{Abbott 2.4.2}
{Abbott}
{subsequence extraction}
{Show that every subsequence of a convergent sequence converges
to the same limit.}

\exstub
{Abbott 2.4.3}
{Abbott}
{limit uniqueness}
{Prove that if a sequence has two convergent subsequences with
different limits then the sequence cannot converge.}

\exstub
{Abbott 2.4.4}
{Abbott}
{subsequence extraction}
{Give an example of a bounded sequence with two subsequences
converging to different limits.}

\exstub
{Abbott 2.4.5}
{Abbott}
{Bolzano–Weierstrass}
{Prove the Bolzano–Weierstrass theorem for sequences in $\mathbb{R}$.}

\exstub
{Abbott 2.4.6}
{Abbott}
{monotone subsequence extraction}
{Show that every sequence has a monotone subsequence.}

\exstub
{Abbott 2.4.7}
{Abbott}
{contradiction}
{Show that if every subsequence has a further subsequence
converging to $L$, then the sequence converges to $L$.}

\exstub
{Abbott 2.4.8}
{Abbott}
{subsequence extraction}
{Construct a bounded sequence with no convergent subsequence
in $\mathbb{Q}$.}

% =========================================================
% BARTLE
% =========================================================

\section*{Exercises from Bartle \& Sherbert}

\exstub
{Bartle 3.5.1}
{Bartle \& Sherbert, \textit{Introduction to Real Analysis}}
{unpack-definition}
{Define a subsequence formally.}

\exstub
{Bartle 3.5.2}
{Bartle \& Sherbert}
{subsequence extraction}
{Show that a subsequence of a convergent sequence converges
to the same limit.}

\exstub
{Bartle 3.5.3}
{Bartle \& Sherbert}
{Bolzano–Weierstrass}
{Prove every bounded sequence in $\mathbb{R}$ has a convergent
subsequence.}

\exstub
{Bartle 3.5.4}
{Bartle \& Sherbert}
{subsequence extraction}
{Construct a sequence with subsequences converging to
different limits.}

\exstub
{Bartle 3.5.5}
{Bartle \& Sherbert}
{limit uniqueness}
{Show that a sequence converges iff every subsequence
converges to the same limit.}

\exstub
{Bartle 3.5.6}
{Bartle \& Sherbert}
{monotone subsequence extraction}
{Show that every sequence has a monotone subsequence.}

\exstub
{Bartle 3.5.7}
{Bartle \& Sherbert}
{Bolzano–Weierstrass}
{Use monotone subsequences to prove Bolzano–Weierstrass.}

\exstub
{Bartle 3.5.8}
{Bartle \& Sherbert}
{tail argument}
{Show that removing finitely many terms does not change
the subsequential limits.}

% =========================================================
% BRUCKNER
% =========================================================

\section*{Exercises from Bruckner}

\exstub
{Bruckner 2.3.1}
{Bruckner, \textit{Real Analysis}}
{subsequence extraction}
{Give an example of a sequence with infinitely many
subsequential limits.}

\exstub
{Bruckner 2.3.2}
{Bruckner}
{subsequence extraction}
{Construct a subsequence converging to the limsup.}

\exstub
{Bruckner 2.3.3}
{Bruckner}
{subsequence extraction}
{Construct a subsequence converging to the liminf.}

\exstub
{Bruckner 2.3.4}
{Bruckner}
{Bolzano–Weierstrass}
{Show that every bounded sequence has a convergent subsequence.}

\exstub
{Bruckner 2.3.5}
{Bruckner}
{subsequence extraction}
{Give a bounded sequence whose subsequential limits form
an interval.}

\exstub
{Bruckner 2.3.6}
{Bruckner}
{limit uniqueness}
{Show that if all subsequences converge to the same limit
then the sequence converges.}

\exstub
{Bruckner 2.3.7}
{Bruckner}
{squeeze argument}
{Use subsequences to prove the squeeze theorem.}

\exstub
{Bruckner 2.3.8}
{Bruckner}
{subsequence extraction}
{Construct subsequences approaching two different limits.}

% =========================================================
% TAO
% =========================================================

\section*{Exercises from Tao}

\exstub
{Tao 6.5.1}
{Tao, \textit{Analysis I}}
{subsequence extraction}
{Define a subsequence formally using strictly increasing indices.}

\exstub
{Tao 6.5.2}
{Tao}
{limit uniqueness}
{Show that a convergent sequence has all subsequences converging
to the same limit.}

\exstub
{Tao 6.5.3}
{Tao}
{Bolzano–Weierstrass}
{Prove every bounded sequence has a convergent subsequence.}

\exstub
{Tao 6.5.4}
{Tao}
{subsequence extraction}
{Construct subsequences converging to limsup and liminf.}

\exstub
{Tao 6.5.5}
{Tao}
{diagonal argument}
{Use a diagonal argument to extract a subsequence satisfying
multiple convergence conditions.}

\exstub
{Tao 6.5.6}
{Tao}
{subsequence extraction}
{Show that a bounded sequence in $\mathbb{R}^n$
has a convergent subsequence.}

\exstub
{Tao 6.5.7}
{Tao}
{compactness argument}
{Use Bolzano–Weierstrass to prove that closed intervals are
sequentially compact.}

\exstub
{Tao 6.5.8}
{Tao}
{diagonal argument}
{Extract a subsequence converging coordinatewise.}

% =========================================================
% LEBL
% =========================================================

\section*{Exercises from Lebl}

\exstub
{Lebl 2.3.1}
{Lebl, \textit{Basic Analysis I}}
{subsequence extraction}
{Define subsequences and give examples.}

\exstub
{Lebl 2.3.2}
{Lebl}
{limit uniqueness}
{Show that subsequences of convergent sequences converge
to the same limit.}

\exstub
{Lebl 2.3.3}
{Lebl}
{subsequence extraction}
{Construct a subsequence converging to a specified limit.}

\exstub
{Lebl 2.3.4}
{Lebl}
{Bolzano–Weierstrass}
{Prove the Bolzano–Weierstrass theorem.}

\exstub
{Lebl 2.3.5}
{Lebl}
{subsequence extraction}
{Show that every bounded sequence in $\mathbb{R}^n$
has a convergent subsequence.}

\exstub
{Lebl 2.3.6}
{Lebl}
{monotone subsequence extraction}
{Show that every sequence contains a monotone subsequence.}

\exstub
{Lebl 2.3.7}
{Lebl}
{squeeze argument}
{Use subsequences to prove the squeeze theorem.}

\exstub
{Lebl 2.3.8}
{Lebl}
{compactness argument}
{Use Bolzano–Weierstrass to prove compactness of
closed intervals.}
